\begin{theorem}[Uniqueness of Peano Systems up to Isomorphism]
\label{thm:uniqueness-of-peano-systems-up-to-isomorphism}
Let $(P,S_P,1_P)$ and $(Q,S_Q,1_Q)$ be Peano systems. Then there exists a
unique bijection $\theta:P\to Q$ such that
\[
\theta(1_P)=1_Q
\]
and
\[
\forall n\in P\;(\theta(S_P(n))=S_Q(\theta(n))).
\]
Thus any two Peano systems are isomorphic.
\hyperref[prf:uniqueness-of-peano-systems-up-to-isomorphism]{Go to proof.}
\end{theorem}
\end{theorembox}

\begin{remark*}[Standard quantified statement]
\[
\exists! \theta:P\to Q\;
\bigl(\theta \text{ is bijective}\land \theta(1_P)=1_Q
\land \forall n\in P\;(\theta(S_P(n))=S_Q(\theta(n)))\bigr).
\]
\end{remark*}

\begin{remark*}[Predicate reading]
Let $A=(P,S_P,1_P)$ and $B=(Q,S_Q,1_Q)$. The theorem asserts
\[
\operatorname{Isomorphism}(A,B),
\]
with unique witness $\theta:P\to Q$ satisfying the displayed base and
successor-preservation clauses.
\end{remark*}

\begin{remark*}[Negated quantified statement]
\[
\begin{aligned}
&\neg\exists! \theta:P\to Q\;
\bigl(\theta \text{ is bijective}\land \theta(1_P)=1_Q\\
&\qquad\qquad\qquad\land
\forall n\in P\;(\theta(S_P(n))=S_Q(\theta(n)))\bigr).
\end{aligned}
\]
\end{remark*}

\begin{remark*}[Failure modes]
\begin{description}
  \item[Exposition.]
  Categoricity fails if no structure-preserving bijection exists between the
  two Peano systems or if more than one such bijection exists.
  \item[Isomorphism failure.]
  \[
  \neg\operatorname{Isomorphism}((P,S_P,1_P),(Q,S_Q,1_Q)).
  \]
  \[
  \begin{aligned}
  &\neg\exists! \theta:P\to Q\;
  \bigl(\theta \text{ is bijective}\land \theta(1_P)=1_Q\\
  &\qquad\qquad\qquad\land
  \forall n\in P\;(\theta(S_P(n))=S_Q(\theta(n)))\bigr).
  \end{aligned}
  \]
\end{description}
\end{remark*}

\begin{remark*}[Interpretation]
All Peano systems have the same structure up to a unique isomorphism. The
elements may be named differently, but the distinguished first element and
successor structure determine the translation.
\end{remark*}

\begin{remark*}[Historical note]
This is the categoricity theorem for Peano systems, following Feferman's
Theorem 3.5 and Simpson's Peano-system presentation
\citep{FefermanNumberSystems1964,simpson2011peanosystems}.
\end{remark*}

\begin{dependencies}
\begin{itemize}
  \item \hyperref[def:peano-system]{Peano System}
  \item \hyperref[thm:general-recursion-theorem-for-peano-system]{General Recursion Theorem}
\end{itemize}
\end{dependencies}
\end{topicbox}
